\section{Introduction}\label{sec:introduction}
Encryption conceals content but need not conceal execution. LLM services produce token streams, cache hit/miss timing and shared-resource activity; agents add destinations, tool exchanges and website actions. A receiving website, a passive network observer and an active API client see different projections \cite{weiss2024prompt,promptpeek2025,earlybird2025,agentprint2025,uitraces2026,jeong2026wra}. The practical question is which deployment decision a favorable defense result actually supports.

A designer may find relevant results without finding an answer to that decision. NetEcho studies downstream re-emission, Dangol measures transport cover, AHB measures executed mobile humanization, and UI Traces studies classifiers fitted to delayed traces \cite{netecho2025,dangol2026workflow,zhu2026turing,uitraces2026}. Each supplies useful evidence. Their percentages and overheads answer different questions, however: what reaches the receiver, what executes, what task outcome remains, and what component the cost measures. A whole-paper defense label obscures these differences.

\paragraph{Central contribution: result-bound deployment judgments}
We develop three complete cases around choosing the final protected interface, deciding whether cover may execute, and selecting an applicable adaptation result and price. Each names the decision, relevant source results, the judgment after condition binding, and the evidence that permits transfer. The resulting rule is to locate the final target-bearing emitter, trace cover to its execution boundary, and bind inference, objective utility and applicable cost to that condition. It preserves positive local conclusions rather than simply declaring the evidence insufficient.

The individual phenomena are already established by the original papers. Our knowledge product is their cross-paper decision mapping: when to test residual features, when to capture a final hop, and when to measure executed task/policy effects. Verified discard before semantic dispatch removes the direct fake-touch mechanism while leaving possible resource and delay effects. The same tested receiver view permits a local adaptation conclusion; a new view needs a bridge. Component costs remain valid for the component, rather than becoming a full-service ranking. These are conditional permissions, not a new composition theorem or alleged corrections of author errors.

\paragraph{Corpus and outputs}
The October 8, 2026 coded collection contains \FinalN{} families: \FinalCore{} core (\FinalA{} workflow-origin and \FinalB{} inherited-runtime/application) plus \FinalAdjacent{} adjacent. Four core components concern privileged diagnostics. Identity/provenance attribution and human--agent distinction are explicitly separated from confidentiality leakage. One result record unifies the method; SCOPE indexes its target, channel, observer, pipeline and inference goal. The outputs are a source--receiver--intervention matrix, three deployment judgments, an applicable-cost subset, and source-deletion/NR sensitivity. The accompanying evidence package links every family, result, source locator and supported claim.

A compact public-trace worked example illustrates a fixed-view prediction decision under archive-receiver ordering. Its task-matched control is inconclusive and its ordinary-site capture bridge is unverified. The separate workflow test favors the complete structured checklist on its original joint endpoint; it establishes no SCOPE review advantage. Failed pilot and retail outcomes are preserved. Remaining acquisition gaps and absent independent researcher adjudication materially limit the synthesis; the paper does not use the corpus count or AI agreement as substitutes for that evidence.

\begin{figure*}[t]
\centering\begin{tikzpicture}[font=\scriptsize,>=Latex]
\tikzset{box/.style={draw,rounded corners,align=center,text width=4.65cm,minimum height=8mm,fill=blue!6},def/.style={box,fill=orange!12},test/.style={box,fill=green!8}}
\foreach \x/\title in {0/{(a) Preserved action structure},5.55/{(b) Downstream re-emission},11.1/{(c) Changed task state}}{
\node[font=\small,anchor=south] at (\x,0.55) {\title};}
\node[box] (a0) at (0,0) {Origin: agent action policy};
\node[def] (a1) at (0,-1.05) {Defense: saved timestamp delays\\action types/order unchanged};
\node[box] (a2) at (0,-2.1) {Receiver: archive / own-site scope\\visible: actions, order, delayed times};
\node[test] (a3) at (0,-3.15) {Test: residual-feature prediction\\live utility remains a separate test};
\node[box] (b0) at (5.55,0) {Origin: inference token stream};
\node[def] (b1) at (5.55,-1.05) {Defense: upstream padding/batching\\application consumes decoded stream};
\node[box] (b2) at (5.55,-2.1) {Receiver: downstream network observer\\visible: group counts / lengths};
\node[test] (b3) at (5.55,-3.15) {Test: paired upstream/final-hop views\\not original timestamp recovery};
\node[box] (c0) at (11.1,0) {Origin: mobile GUI agent actions};
\node[def] (c1) at (11.1,-1.05) {Defense: executed fake touches\\can trigger an unintended transition};
\node[box] (c2) at (11.1,-2.1) {Receiver: app touch-event detector\\visible: events and interaction timing};
\node[test] (c3) at (11.1,-3.15) {Test: detection and terminal goals\\retain failures / transition changes};
\foreach \p in {a,b,c}{\draw[->] (\p0)--(\p1);\draw[->] (\p1)--(\p2);\draw[->] (\p2)--(\p3);}
\end{tikzpicture}

\caption{Three source-bound situations: preserved actions (UI Traces), downstream group/length re-emission (NetEcho), and state-changing executed touches (AHB). Signal path and task-state effect determine different deployment tests.}
\label{fig:system-model}
\end{figure*}
